\section{Template}

\begin{wrapfigure}{r}{0.5\textwidth}
  \begin{center}
    \includegraphics[width=0.45\textwidth]{img/logo/EPFL_logo.png}
  \end{center}
  \caption{Caption}
\label{fig:1}
\end{wrapfigure}

\begin{equation}
	\vec{B} = \mu_0 \mu_{r}\vec{H}
	\label{eqn:1}
\end{equation}

\begin{figure}[htpb]
	\centering
	\includegraphics[width=0.8\textwidth]{img/logo/EPFL_logo.png}
	\caption{Montage de l'expérience \cite{ref2}.}
	\label{fig:montage}
\end{figure}

\begin{table}[htbp]
    \centering
    \sisetup{
        separate-uncertainty = true,
        round-mode = uncertainty,
        round-precision = 1
    }
    \begin{tabular}{
        S[table-format=1.0]
        S[table-format=1.1(1)]
        S[table-format=3.0(1)]
    }
        \toprule
        {Isolants} & {$B_{r}$ [\si{\milli\tesla}]} & {$H_{c}$ [\si{\ampere\per\meter}]} \\
        \midrule
        1 & 3.57(23) & 271(38) \\
        2 & 1.69(28) & 209(41) \\
        4 & 0.76(21) & 158(54) \\
        8 & 0.30(20) & 102(79) \\
        \bottomrule
    \end{tabular}
    \caption{Champ rémanent $B_r$ et champ coercitif $H_c$ en fonction du nombre
    d'isolants}
    \label{tab:champs_max}
\end{table}
